\documentclass[12pt, titlepage]{article}

\usepackage{booktabs}
\usepackage{tabularx}
\usepackage{hyperref}
\hypersetup{
    colorlinks,
    citecolor=blue,
    filecolor=black,
    linkcolor=red,
    urlcolor=blue
}
\usepackage[round]{natbib}

\input{../Comments.text}
\input{../Common.text}

\begin{document}

\title{System Verification and Validation Plan for \progname{}} 
\author{\authname}
\date{\today}
	
\maketitle

\pagenumbering{roman}

\section*{Revision History}

\begin{tabularx}{\textwidth}{p{3cm}p{2cm}X}
\toprule {\bf Date} & {\bf Version} & {\bf Notes}\\
\midrule
Date 1 & 1.0 & Notes\\
Date 2 & 1.1 & Notes\\
\bottomrule
\end{tabularx}

~\\
\wss{The intention of the VnV plan is to increase confidence in the software.
However, this does not mean listing every verification and validation technique
that has ever been devised.  The VnV plan should also be a \textbf{feasible}
plan. Execution of the plan should be possible with the time and team available.
If the full plan cannot be completed during the time available, it can either be
modified to ``fake it'', or a better solution is to add a section describing
what work has been completed and what work is still planned for the future.}

\wss{The VnV plan is typically started after the requirements stage, but before
the design stage.  This means that the sections related to unit testing cannot
initially be completed.  The sections will be filled in after the design stage
is complete.  the final version of the VnV plan should have all sections filled
in.}

\newpage

\tableofcontents

\listoftables
\wss{Remove this section if it isn't needed}

\listoffigures
\wss{Remove this section if it isn't needed}

\newpage

\section{Symbols, Abbreviations, and Acronyms}

\renewcommand{\arraystretch}{1.2}
\begin{tabular}{l l} 
  \toprule		
  \textbf{symbol} & \textbf{description}\\
  \midrule 
  T & Test\\
  \bottomrule
\end{tabular}\\

\wss{symbols, abbreviations, or acronyms --- you can simply reference the SRS
  \citep{SRS} tables, if appropriate}

\wss{Remove this section if it isn't needed}

\newpage

\pagenumbering{arabic}

This document ... \wss{provide an introductory blurb and roadmap of the
  Verification and Validation plan}

\section{General Information}

\subsection{Summary}

\wss{Say what software is being tested.  Give its name and a brief overview of
  its general functions.}

\subsection{Objectives}

\wss{State what is intended to be accomplished.  The objective will be around
  the qualities that are most important for your project.  You might have
  something like: ``build confidence in the software correctness,''
  ``demonstrate adequate usability.'' etc.  You won't list all of the qualities,
  just those that are most important.}

\wss{You should also list the objectives that are out of scope.  You don't have 
the resources to do everything, so what will you be leaving out.  For instance, 
if you are not going to verify the quality of usability, state this.  It is also 
worthwhile to justify why the objectives are left out.}

\wss{The objectives are important because they highlight that you are aware of 
limitations in your resources for verification and validation.  You can't do everything, 
so what are you going to prioritize?  As an example, if your system depends on an 
external library, you can explicitly state that you will assume that external library 
has already been verified by its implementation team.}

\subsection{Custom Documents}

\wss{Summarize the CDs (Custom Documents) that were tackled by this project.  CDs 
can include usability testing, code walkthroughs, user documentation, formal
proof, GenderMag personas, Design Thinking, etc.  CDs should have already
been approved by the course instructor as included in your problem statement.
You can use a pull request to update your CDs (in TeamComposition.csv or
Repos.csv) if your plan changes as a result of the VnV planning exercise.}

\subsection{Relevant Documentation}

\wss{Reference relevant documentation.  This will definitely include your SRS
  and your other project documents (design documents, like MG, MIS, etc).  You
  can include these even before they are written, since by the time the project
  is done, they will be written.  You can create BibTeX entries for your
  documents and within those entries include a hyperlink to the documents.}

\citet{SRS}

\wss{Don't just list the other documents.  You should explain why they are relevant and 
how they relate to your VnV efforts.}

\section{Plan}

\wss{Introduce this section.  You can provide a roadmap of the sections to
  come.}

\subsection{Verification and Validation Team}

\wss{Your teammates.  Maybe your supervisor.
  You should do more than list names.  You should say what each person's role is
  for the project's verification.  A table is a good way to summarize this information.}

\subsection{SRS Verification}

\wss{List any approaches you intend to use for SRS verification.  This may
  include ad hoc feedback from reviewers, like your classmates (like your
  primary reviewer), or you may plan for something more rigorous/systematic.}

\wss{If you have a supervisor for the project, you shouldn't just say they will
read over the SRS.  You should explain your structured approach to the review.
Will you have a meeting?  What will you present?  What questions will you ask?
Will you give them instructions for a task-based inspection?  Will you use your
issue tracker?}

\wss{Maybe create an SRS checklist?}

\subsection{Design Verification}

\wss{Plans for design verification}

\wss{The review will include reviews by your classmates}

\wss{Create a checklists?}

\subsection{Verification and Validation Plan Verification}

\wss{The verification and validation plan is an artifact that should also be
verified.  Techniques for this include review and mutation testing.}

\wss{The review will include reviews by your classmates}

\wss{Create a checklists?}

\subsection{Implementation Verification}

\wss{You should at least point to the tests listed in this document and the unit
  testing plan.}

\wss{In this section you would also give any details of any plans for static
  verification of the implementation.  Potential techniques include code
  walkthroughs, code inspection, static analyzers, etc.}

\wss{The final class presentation in CAS 741 could be used as a code
walkthrough.  There is also a possibility of using the final presentation (in
CAS741) for a partial usability survey.}

\subsection{Automated Testing and Verification Tools}

\wss{What tools are you using for automated testing.  Likely a unit testing
  framework and maybe a profiling tool, like ValGrind.  Other possible tools
  include a static analyzer, make, continuous integration tools, test coverage
  tools, etc.  Explain your plans for summarizing code coverage metrics.
  Linters are another important class of tools.  For the programming language
  you select, you should look at the available linters.  There may also be tools
  that verify that coding standards have been respected, like flake9 for
  Python.}

\wss{If you have already done this in the development plan, you can point to
that document.}

\wss{The details of this section will likely evolve as you get closer to the
  implementation.}

\subsection{Software Validation}

\wss{If there is any external data that can be used for validation, you should
  point to it here.  If there are no plans for validation, you should state that
  here.}

\wss{You might want to use review sessions with the stakeholder to check that
the requirements document captures the right requirements.  Maybe task based
inspection?}

\wss{For those capstone teams with an external supervisor, the Rev 0 demo should 
be used as an opportunity to validate the requirements.  You should plan on 
demonstrating your project to your supervisor shortly after the scheduled Rev 0 demo.  
The feedback from your supervisor will be very useful for improving your project.}

\wss{For teams without an external supervisor, user testing can serve the same purpose 
as a Rev 0 demo for the supervisor.}

\wss{This section might reference back to the SRS verification section.}

\section{System Tests}

\wss{There should be text between all headings, even if it is just a roadmap of
the contents of the subsections.}

\subsection{Tests for Functional Requirements} \label{SecFRTests}

% \testdef marks where a test is defined; \testref links to it from the
% traceability table in Section~\ref{SecTraceability}.
\newcommand{\testdef}[1]{\hypertarget{test:#1}{#1}}
\newcommand{\testref}[1]{\hyperlink{test:#1}{#1}}

The tests below are categorized by the components defined in Section S.2 of the
SRS \citep{SRS}, so that every functional requirement is covered. Each test is a
system-level test of the deployed system, following a workflow that a club
member performs through the Toolbox Web App, and verifies several related
requirements. The tests also verify the functional safety and security
requirements that the Hazard Analysis Roadmap plans for the capstone.

\subsubsection{Web Server (S.2.1)}

\begin{enumerate}

\item \testdef{T.FR.WS.1.1}\\
\textbf{Control:} Manual\\
\textbf{Initial State:} System deployed. User not logged in.\\
\textbf{Input:} User opens the public landing page, then attempts to open a club
service. User then logs in with valid credentials, uses the service, and logs
out.\\
\textbf{Output:} The landing page is shown without logging in. The club service
redirects to the login page until the user logs in, and is no longer accessible
after logout.\\
\textbf{Derivation:} Verifies FR-WS-1, FR-WS-2, FR-WS-3.

\end{enumerate}

\subsubsection{Access Control System (S.2.2)}

\begin{enumerate}

\item \testdef{T.FR.AC.2.1}\\
\textbf{Control:} Manual\\
\textbf{Initial State:} An executive is logged in.\\
\textbf{Input:} The executive creates an account, first with a password that does
not meet the password requirements, then with one that does. The executive
changes the account's role and team, ends its session from the server side, and
deactivates it.\\
\textbf{Output:} The invalid password is rejected with the reason. The account
always has exactly one role and at most one team. Ending the session logs the
user out, and the deactivated account cannot log in.\\
\textbf{Derivation:} Verifies FR-AC-1, FR-AC-2, FR-AC-4, FR-AC-5, SR-AC-2,
SR-WS-5.

\item \testdef{T.FR.AC.2.2}\\
\textbf{Control:} Automatic\\
\textbf{Initial State:} One account for each role, on two different teams. Each
team has a current bridge model.\\
\textbf{Input:} Each account requests each restricted operation (member
management, rule configuration, season archiving, test result review) and each
team's current model.\\
\textbf{Output:} Each request is allowed only for the roles that the SRS permits.
Refused requests show an error and change nothing. Only members of a team can
access its current models.\\
\textbf{Derivation:} Verifies FR-AC-3, FR-AC-6, FR-BR-2.

\end{enumerate}

\subsubsection{Toolbox Broker (S.2.3)}

\begin{enumerate}

\item \testdef{T.FR.BR.3.1}\\
\textbf{Control:} Manual\\
\textbf{Initial State:} System has just been started. User logged in.\\
\textbf{Input:} User opens each tool. One service is then stopped, and the user
requests that service and another tool.\\
\textbf{Output:} Every tool opens. The stopped service returns an error message
instead of hanging, while the other tool keeps working.\\
\textbf{Derivation:} Verifies FR-BR-1, FR-BR-3, FR-BR-4.

\end{enumerate}

\subsubsection{Knowledge Repository (S.2.4)}

\begin{enumerate}

\item \testdef{T.FR.KR.4.1}\\
\textbf{Control:} Manual\\
\textbf{Initial State:} User logged in. Repository contains items with known
metadata.\\
\textbf{Input:} User uploads files of common types with metadata, and one file
without the required metadata. User then searches and filters the repository,
previews a file, and downloads one.\\
\textbf{Output:} Valid uploads are stored with their metadata, and the upload
without required metadata is rejected with a message. Searches return exactly
the matching items, and downloaded files are identical to the uploaded ones.\\
\textbf{Derivation:} Verifies FR-KR-1, FR-KR-2, FR-KR-3, FR-KR-4, SR-KR-4.

\item \testdef{T.FR.KR.4.2}\\
\textbf{Control:} Manual\\
\textbf{Initial State:} A member has uploaded an item.\\
\textbf{Input:} The uploader edits the item, another member attempts to edit it,
and an executive archives it. The repository's storage is then made unavailable,
and a member uploads a file.\\
\textbf{Output:} Each change asks for confirmation first, and the other member's
edit is refused. The archived item no longer appears in search results but can
still be opened. The failed upload shows an error message.\\
\textbf{Derivation:} Verifies FR-KR-5, SR-KR-2, SR-KR-5.

\end{enumerate}

\subsubsection{Model Analysis System (S.2.5)}

\begin{enumerate}

\item \testdef{T.FR.MA.5.1}\\
\textbf{Control:} Manual\\
\textbf{Initial State:} Team member logged in.\\
\textbf{Input:} User uploads invalid and corrupted model files, then a reference
STEP model with known properties, confirms the import summary, and rotates,
zooms, and pans the 3D view.\\
\textbf{Output:} Invalid files are rejected with the reason. The import summary
shows extracted properties that match the reference model's known values before
any analysis is run. The 3D view responds to each interaction and matches the
model in CAD. A model that exceeds the processing time limit is stopped with a
message.\\
\textbf{Derivation:} Verifies FR-MA-1, FR-MA-2, FR-MA-3, SR-MA-1, SR-MA-3.

\item \testdef{T.FR.MA.5.2}\\
\textbf{Control:} Manual\\
\textbf{Initial State:} Two teams exist, with no current models.\\
\textbf{Input:} A team member uploads a current model. A teammate, a member of the
other team, and a former team member try to open it. An executive then archives
the season's models, and a general member browses past designs.\\
\textbf{Output:} Only the teammate can open the current model. After archiving, it
becomes a past design recording its year and team, which every member can view
and filter but not modify.\\
\textbf{Derivation:} Verifies FR-MA-4, FR-MA-5, FR-MA-6, FR-AC-6.

\end{enumerate}

\subsubsection{Performance Analysis System (S.2.5.1)}

\begin{enumerate}

\item \testdef{T.FR.PA.6.1}\\
\textbf{Control:} Manual\\
\textbf{Initial State:} A reference model whose failure load can be calculated
by hand is stored, with two material data sets.\\
\textbf{Input:} User runs a performance analysis of the reference model with
each material set, then with the calculation service unavailable.\\
\textbf{Output:} The failure load and load-to-mass ratio match the values
calculated by hand using SRS S.7.2, and change with the selected material set.
The assumptions and limitations are shown with the results. When the service is
unavailable, an error is shown instead of an estimate.\\
\textbf{Derivation:} Verifies FR-PA-1, FR-PA-2, FR-PA-3, FR-PA-4, FR-MT-4,
SR-PA-1.

\end{enumerate}

\subsubsection{Rule Compliance System (S.2.5.2)}

\begin{enumerate}

\item \testdef{T.FR.RC.7.1}\\
\textbf{Control:} Manual\\
\textbf{Initial State:} Executive logged in. Test models stored that pass, violate,
and sit exactly at the competition rule limits.\\
\textbf{Input:} The executive imports and edits a rule set for a competition
year. A member then runs a compliance check on each test model.\\
\textbf{Output:} The rule set is stored under its year with the edits applied.
Each rule passes or fails as given by SRS S.7.3, and each failure shows the
measured and allowed values. Rules that cannot be checked from the model are
listed as manual checks.\\
\textbf{Derivation:} Verifies FR-RC-1, FR-RC-2, FR-RC-3, FR-RC-4, FR-RC-5.

\end{enumerate}

\subsubsection{Material Testing Manager (S.2.6)}

\begin{enumerate}

\item \testdef{T.FR.MT.8.1}\\
\textbf{Control:} Manual\\
\textbf{Initial State:} Material test results from a previous year are stored.\\
\textbf{Input:} User enters test results through the form, and imports a file
that includes invalid rows. User views the comparison with the previous year,
and an executive flags one result as invalid.\\
\textbf{Output:} Valid results are stored, and wood results require a grain
orientation. Invalid rows are rejected and identified. The comparison statistics
match values calculated by hand, and the material properties are recomputed
without the flagged result, which remains stored.\\
\textbf{Derivation:} Verifies FR-MT-1, FR-MT-2, FR-MT-3, FR-MT-4, FR-MT-5,
SR-MT-2.

\end{enumerate}

\subsubsection{Knowledge Transfer Manager (S.2.7)}

\begin{enumerate}

\item \testdef{T.FR.KT.9.1}\\
\textbf{Control:} Manual\\
\textbf{Initial State:} Executive logged in.\\
\textbf{Input:} The executive opens their exit document template, saves a draft,
and submits it, first with a reference to an item that does not exist, then with
valid references. Another executive views the outstanding exit documents.\\
\textbf{Output:} A template exists for each position, and the draft keeps its
content. The invalid reference is rejected. The valid submission is published
only once the repository confirms it was saved, and it can then be found by
search. The outstanding list no longer includes the executive.\\
\textbf{Derivation:} Verifies FR-KT-1, FR-KT-2, FR-KT-4, FR-KT-5, SR-KT-2,
SR-KT-5.

\item \testdef{T.FR.KT.9.2}\\
\textbf{Control:} Manual\\
\textbf{Initial State:} An executive holds a position with a role record.\\
\textbf{Input:} The executive and another member each try to update the role
record. The executive's account is then archived, first without a successor,
then after a successor is assigned.\\
\textbf{Output:} Only the executive's update is saved. Archiving is blocked until a
successor is assigned, after which the successor can update the role record.\\
\textbf{Derivation:} Verifies FR-KT-3, FR-KT-4, SR-KT-3.

\end{enumerate}

\subsection{Tests for Nonfunctional Requirements}

\wss{The nonfunctional requirements for accuracy will likely just reference the
  appropriate functional tests from above.  The test cases should mention
  reporting the relative error for these tests.  Not all projects will
  necessarily have nonfunctional requirements related to accuracy.}

\wss{For some nonfunctional tests, you won't be setting a target threshold for
passing the test, but rather describing the experiment you will do to measure
the quality for different inputs.  For instance, you could measure speed versus
the problem size.  The output of the test isn't pass/fail, but rather a summary
table or graph.}

\wss{Tests related to usability could include conducting a usability test and
  survey.  The survey will be in the Appendix.}

\wss{Static tests, review, inspections, and walkthroughs, will not follow the
format for the tests given below.}

\wss{If you introduce static tests in your plan, you need to provide details.
How will they be done?  In cases like code (or document) walkthroughs, who will
be involved? Be specific.}

\subsubsection{Area of Testing1}
		
\paragraph{Title for Test}

\begin{enumerate}

\item{test-id1\\}

Type: Functional, Dynamic, Manual, Static etc.
					
Initial State: 
					
Input/Condition: 
					
Output/Result: 
					
How test will be performed: 
					
\item{test-id2\\}

Type: Functional, Dynamic, Manual, Static etc.
					
Initial State: 
					
Input: 
					
Output: 
					
How test will be performed: 

\end{enumerate}

\subsubsection{Area of Testing2}

...

\subsection{Traceability Between Test Cases and Requirements} \label{SecTraceability}

\begin{table}[htbp]
\centering
\caption{Functional Requirement Traceability Matrix}
\label{TblTraceFR}
\begin{tabular}{|l|p{0.65\textwidth}|}
\hline
\textbf{Test Case ID} & \textbf{Requirements Verified}\\
\hline
\testref{T.FR.WS.1.1} & FR-WS-1, FR-WS-2, FR-WS-3\\
\hline
\testref{T.FR.AC.2.1} & FR-AC-1, FR-AC-2, FR-AC-4, FR-AC-5, SR-AC-2, SR-WS-5\\
\hline
\testref{T.FR.AC.2.2} & FR-AC-3, FR-AC-6, FR-BR-2\\
\hline
\testref{T.FR.BR.3.1} & FR-BR-1, FR-BR-3, FR-BR-4\\
\hline
\testref{T.FR.KR.4.1} & FR-KR-1, FR-KR-2, FR-KR-3, FR-KR-4, SR-KR-4\\
\hline
\testref{T.FR.KR.4.2} & FR-KR-5, SR-KR-2, SR-KR-5\\
\hline
\testref{T.FR.MA.5.1} & FR-MA-1, FR-MA-2, FR-MA-3, SR-MA-1, SR-MA-3\\
\hline
\testref{T.FR.MA.5.2} & FR-MA-4, FR-MA-5, FR-MA-6, FR-AC-6\\
\hline
\testref{T.FR.PA.6.1} & FR-PA-1, FR-PA-2, FR-PA-3, FR-PA-4, FR-MT-4, SR-PA-1\\
\hline
\testref{T.FR.RC.7.1} & FR-RC-1, FR-RC-2, FR-RC-3, FR-RC-4, FR-RC-5\\
\hline
\testref{T.FR.MT.8.1} & FR-MT-1, FR-MT-2, FR-MT-3, FR-MT-4, FR-MT-5, SR-MT-2\\
\hline
\testref{T.FR.KT.9.1} & FR-KT-1, FR-KT-2, FR-KT-4, FR-KT-5, SR-KT-2, SR-KT-5\\
\hline
\testref{T.FR.KT.9.2} & FR-KT-3, FR-KT-4, SR-KT-3\\
\hline
\end{tabular}
\end{table}

The safety and security requirements that the Hazard Analysis Roadmap delays
until after the capstone are not tested: SR-KR-3, SR-MA-2, SR-RC-1, SR-KT-1,
SR-KT-4, SR-MT-1, and SR-MT-3.

% Non-functional requirement traceability (Section 4.2 tests) to be added below.

\section{Unit Test Description}

\wss{This section should not be filled in until after the MIS (detailed design
  document) has been completed.}

\wss{Reference your MIS (detailed design document) and explain your overall
philosophy for test case selection.}  

\wss{To save space and time, it may be an option to provide less detail in this section.  
For the unit tests you can potentially layout your testing strategy here.  That is, you 
can explain how tests will be selected for each module.  For instance, your test building 
approach could be test cases for each access program, including one test for normal behaviour 
and as many tests as needed for edge cases.  Rather than create the details of the input 
and output here, you could point to the unit testing code.  For this to work, you code 
needs to be well-documented, with meaningful names for all of the tests.}

\subsection{Unit Testing Scope}

\wss{What modules are outside of the scope.  If there are modules that are
  developed by someone else, then you would say here if you aren't planning on
  verifying them.  There may also be modules that are part of your software, but
  have a lower priority for verification than others.  If this is the case,
  explain your rationale for the ranking of module importance.}

\subsection{Tests for Functional Requirements}

\wss{Most of the verification will be through automated unit testing.  If
  appropriate specific modules can be verified by a non-testing based
  technique.  That can also be documented in this section.}

\subsubsection{Module 1}

\wss{Include a blurb here to explain why the subsections below cover the module.
  References to the MIS would be good.  You will want tests from a black box
  perspective and from a white box perspective.  Explain to the reader how the
  tests were selected.}

\begin{enumerate}

\item{test-id1\\}

Type: \wss{Functional, Dynamic, Manual, Automatic, Static etc. Most will
  be automatic}
					
Initial State: 
					
Input: 
					
Output: \wss{The expected result for the given inputs}

Test Case Derivation: \wss{Justify the expected value given in the Output field}

How test will be performed: 
					
\item{test-id2\\}

Type: \wss{Functional, Dynamic, Manual, Automatic, Static etc. Most will
  be automatic}
					
Initial State: 
					
Input: 
					
Output: \wss{The expected result for the given inputs}

Test Case Derivation: \wss{Justify the expected value given in the Output field}

How test will be performed: 

\item{...\\}
    
\end{enumerate}

\subsubsection{Module 2}

...

\subsection{Tests for Nonfunctional Requirements}

\wss{If there is a module that needs to be independently assessed for
  performance, those test cases can go here.  In some projects, planning for
  nonfunctional tests of units will not be that relevant.}

\wss{These tests may involve collecting performance data from previously
  mentioned functional tests.}

\subsubsection{Module ?}
		
\begin{enumerate}

\item{test-id1\\}

Type: \wss{Functional, Dynamic, Manual, Automatic, Static etc. Most will
  be automatic}
					
Initial State: 
					
Input/Condition: 
					
Output/Result: 
					
How test will be performed: 
					
\item{test-id2\\}

Type: Functional, Dynamic, Manual, Static etc.
					
Initial State: 
					
Input: 
					
Output: 
					
How test will be performed: 

\end{enumerate}

\subsubsection{Module ?}

...

\subsection{Traceability Between Test Cases and Modules}

\wss{Provide evidence that all modules have been considered.}
				
\bibliographystyle{plainnat}

\bibliography{../../refs/References}

\newpage

\section{Appendix}

This is where you can place additional information.

\subsection{Symbolic Parameters}

The definition of the test cases will call for SYMBOLIC\_CONSTANTS.
Their values are defined in this section for easy maintenance.

\subsection{Usability Survey Questions?}

\wss{This is a section that would be appropriate for some projects.}

\newpage{}

\section{Appendix --- Generative AI Use Disclosure}

\wss{ChatGPT (version 5) was used to brainstorm ideas for the template for this disclose section.}

\subsection{AI Tools Used}

\wss{Give the name of the AI tool(s) used and the version.}

\subsection{How and Where Used}

\wss{Explain what the tool(s) were used for.  Examples include: brainstorming,
explaining a technical concept, reviewing the document, generating or modifying
text, generating or modifying diagrams, identifying errors or omissions or
inconsistencies.}

\wss{For each of the uses, identify which parts of the document were affected.}

\wss{You do not need to report every interaction, but you should disclose the 
uses that materially contributed to your work}

\subsection{Verification of AI-Generated Content}

\wss{\begin{itemize}
    \item Describe how the team checked AI-generated or AI-assisted material
    before including it in the document.
    \item In particular, verify factual claims, technical assertions,
    requirements, calculations, references, and other information for which an
    AI system may produce incorrect or fabricated results.
    \item \textbf{The team is responsible for the accuracy, completeness,
consistency, and appropriateness of the submitted document, regardless of
whether material was generated by a human or an AI system.}
\end{itemize}}

\newpage{}

\section{Appendix --- Reflection}

\wss{This section is not required for CAS 741}

The information in this section will be used to evaluate the team members on the
graduate attribute of Lifelong Learning.

\input{../Reflection.text}

\begin{enumerate}
  \item What went well while writing this deliverable? 
  \item What pain points did you experience during this deliverable, and how
    did you resolve them?
  \item What knowledge and skills will the team collectively need to acquire to
  successfully complete the verification and validation of your project?
  Examples of possible knowledge and skills include dynamic testing knowledge,
  static testing knowledge, specific tool usage, Valgrind etc.  You should look to
  identify at least one item for each team member.
  \item For each of the knowledge areas and skills identified in the previous
  question, what are at least two approaches to acquiring the knowledge or
  mastering the skill?  Of the identified approaches, which will each team
  member pursue, and why did they make this choice?
\end{enumerate}

\end{document}